\section{Computing the two-point angular functions}\label{sec:2pcf}

Two point angular correlations are computed with the \texttt{pycorr} software, based on the \texttt{CORRFUNC}\footnote{\textcolor{blue}{\href{https://github.com/manodeep/Corrfunc}{\texttt{github.com/manodeep/Corrfunc}}}} \cite{Sinha2019corrfunc1, Sinha2019corrfunc2} engine. HSC and DESI's survey footprints are extracted using multi-order coverage maps (MOC) \cite{fernique2015moc} made with the \texttt{mocpy}\footnote{\textcolor{blue}{\href{https://github.com/cds-astro/mocpy}{\texttt{github.com/cds-astro/mocpy}}}} python package. The areas of interest for this study represent the intersection of DESI's MOC with HSC's MOC.
%For HSC, we choose a maximum resolution order of $12$, which gives an approximate angular resolution of $\sim 51$ arcsec, still efficiently covering the footprint and survey mask without over-removing areas. For DESI, we use a maximum order of $10$, since DESI's footprint is already well characterized with the random catalog. => Seems like this is not super relevant, maybe remove or place in a footnote ?
The computations of the two-point angular functions are split up into four patches of sky, corresponding to spatially close areas, in order to speed up computing time. Concretely, this corresponds to the \texttt{XMM} field, the \texttt{HECTOMAP} field, the \texttt{VVDS} field, and the joint field composed of \texttt{GAMA09H}, \texttt{GAMA15H} and \texttt{WIDE12H}. We use $N_\mathrm{b}=32$ logarithmically-spaced angular bins, with edges spanning $0.1$ to $10$ $\hMpc$ , in transverse comoving scales.

The DESI angular auto-correlations are obtained by combining the auto-correlation measured in the north galactic cap (NGC) and south galactic cap (SGC) of the dataset used. Two point correlation functions are computed independently per DESI tracer, and include standard weights and the FKP weights \cite{FKPWeights1994}. 
The weighting used is described in \cite{DESIConstructionLSS2024}: $w_\mathrm{tot}=w_\mathrm{comp}w_\mathrm{sys}w_\mathrm{zfail}w_\mathrm{FKP}$. Here, $w_\mathrm{comp}$ is used for completeness, $w_\mathrm{sys}$ accounts for target density fluctuations due to imaging conditions and $w_\mathrm{zfail}$ encompasses the relative redshift success rate. Given the redshift ranges used in this analysis, the FKP weights $w_{\mathrm{FKP}}$ are not strictly necessary, since their purpose is to describe how the number density changes with redshift, and the redshift bins are fine. It is important to note that fiber assignment incompleteness affects our measurements on these small scales, especially for low completeness tracers such as ELGs. 
This effect is all the more prominent for DR1 ELGs, which is why we decide to not use them in this analysis. 
In short, the number of targets who obtain a fiber on a single DESI tile while observing is different from the total number of possible targets in that target class within the tile. 
To mitigate this effect, one could use Pairwise Inverse Probability (PIP, \cite{BianchiPIP2017}) weights when computing two point correlation functions, but this is only doable when performing auto-correlations.
Since cross-correlations are performed with the HSC catalog, PIP weighting is not available, though it is noteworthy that most of HSC's footprint is already within DESI areas with high or maximal completeness, as shown in figure \ref{fig:footprint}. 
This kind of clustering redshift study will benefit from subsequent data releases of DESI, alleviating further this issue. 
To verify this introduces no apparent bias in the measurements, two scale cuts are included in this work: $0.3-3\hMpc$ and $1-5\hMpc$. 
The fiber collisions become an issue below $\theta_\mathrm{fib}=0.05^\circ$ \cite{bianchi2025characterizationdesifiberassignment, Pinon2025FiberAssign}. 
The second scale cut allows for no leakage of fiber collisions in the scale cut up to $z\sim0.45$, with minor contamination for larger redshifts. 
The measurement consistency shown in section \ref{sec:results} supports that no bias is introduces by using measurements at smaller scales.

With the HSC catalog, the lensing weights are included in our two-point estimators, following the same choices as \cite{HSCClusteringRau2023}. 
Jackknife estimates of the covariance matrix are computed to estimate error bars, including the correction given by \cite{Mohammad2022jackknife}, which is the fiducial approach implemented in \texttt{pycorr}. 
Per patch of the sky, $N_{\mathrm{JK}}=100$ "delete-one" jackknife realizations are performed: in each iteration, a random area of sky in the footprint is removed and the two-point angular statistic is computed with that area excluded. 
Removed areas are determined through a KMeans sub-sampler algorithm splitting up the footprint into $N_{\mathrm{JK}}$ areas at \texttt{HEALPix} \cite{Gorski2005HEALPix} nside resolution of $256$ on each of the patches of sky. The covariance matrix obtained on the associated angular two-point correlation functions can be written as such \cite{Mohammad2022jackknife}:

\begin{equation}
    \mathbf{C}_{ij}=\frac{N_{\mathrm{JK}}-1}{N_{\mathrm{JK}}}\sum^{N_{\mathrm{JK}}}_{n=1}(\omega^{(n)}_i-\langle\omega_i\rangle)(\omega^{(n)}_j-\langle\omega_j\rangle) %^\mathrm{T} (could also use matrix notation instead of i,j, not sure which is better to understand
\end{equation}

where $({N_{\mathrm{JK}}-1})/{N_{\mathrm{JK}}}$ is the fraction of the footprint's area used to compute any jackknife realization, $\omega^{(n)}_i$ is the measurement in angular bin $i\in\llbracket 1,N_b \rrbracket$  over the footprint excluding area $n$ and $\langle\omega_i\rangle$ is the mean value over all measurements in bin $i$ from excluded regions. The Davis-Peebles \cite{DavisPeebles1983} estimator is chosen to compute the two-point angular correlation functions:

\begin{equation}
    \hat{\omega}_{xy}(r):=\frac{\mathrm{D}_x\mathrm{D}_y(r)}{\mathrm{R}_x\mathrm{D}_y(r)}-1
\end{equation}

where $\mathrm{D}_x\mathrm{D}_y$ are normalized data-data pair counts and $\mathrm{D}_x\mathrm{R}_y$ are the normalized data-random pair counts. Although this estimator is not as low bias and variance as the broadly used Landy-Szalay \cite{LandySzalay1993} estimator:

\begin{equation}
    \hat{\omega}^{\mathrm{LS}}_{xy}(r):=\frac{\mathrm{D}_x\mathrm{D}_y(r)-\mathrm{R}_x\mathrm{D}_y(r)-\mathrm{D}_x\mathrm{R}_y(r)}{\mathrm{R}_x\mathrm{R}_y(r)}+1, 
\end{equation}

the Davis-Peebles estimator allows this study to not rely on the random catalog of the photometric sample, and only rely on the random catalog for our spectroscopic tracers. A random catalog for HSC is still used when computing the auto-correlations on smaller tomographic bins $\bo_{pp}^{\mathrm{photo}-z}$, in the context of photometric galaxy bias mitigation (see section  \ref{sec:modeling:photoz_gal_bias}). This random catalog has no redshift assignments, so it is less robust than the DESI catalog, but captures the footprint sufficiently well given our measurement. \footnote{More details about the random points used can be found on the HSC data release 3 portal: \href{https://hsc-release.mtk.nao.ac.jp/doc/index.php/random-points__pdr3/}{\textcolor{blue}{hsc-release.mtk.nao.ac.jp/random-points\_\_pdr3/}}.} Conveniently, the Davis-Peebles estimator is also faster to compute. At least 10 times more randoms from the DESI catalogs than data points are used when computing the two-point correlation functions\footnote{This corresponds to four catalogs of randoms from the Large Scale Structure catalogs of DESI \cite{DESIConstructionLSS2024}. Each catalog corresponds to a number density of $2500$ deg$^{-2}$, so the random catalog number density is $10000$ deg$^{-2}$ before redshift binning.}, and for low-density datasets such as QSOs, this can reach up to 100 times more. %should we give exact numbers here ? => I think it's okay to specify that we use 4 random catalog realizations, as given in the footnote. These data products are public, so the information on the exact numbers are there. The ratio data/random can slightly change per tracer or redshift bin. --> SGG: no need to go to that level of detail. 